\chapter{Sprays, connections and fields of projectors}
\label{chap:connections}

In this chapter we introduce the second-order structures of a manifold $M$ modeled on a locally
convex space $E$: connections, given in charts by their Christoffel symbols, and sprays, given in
charts by their local parts. We define the covariant derivative of vector fields and of lifts of
curves, the torsion and the curvature, and the geodesics of a connection. All the objects are
defined chart by chart, subject to a change of charts law, and the main point in each case is
that the resulting object does not depend on the chart. We follow Lang \cite{Lan99}
(\S\S IV.3, VIII.1--3, IX.1) and Chapter~6 of \cite{Lar22}.

The last section is closer to \cite{Lar22}: a smooth field of projectors onto the tangent spaces
of an embedded submanifold of a locally convex space defines a symmetric connection, whose
covariant derivative is obtained by differentiating in the ambient space and projecting, and
whose geodesics are the curves with normal acceleration.

Two features of the locally convex setting shape the chapter. First, there is no good topology on
spaces of continuous linear operators, so bilinear families such as $x \mapsto \Gamma_x$ or
$p \mapsto E_p$ are encoded as maps of several variables, linear in some of them, with joint
smoothness. Second, there is no existence theorem for ordinary differential equations in general
locally convex spaces; consequently we define geodesics by their equation, but we never claim that
they exist for given initial data, and parallel transport is not discussed.

\section{Connections and the covariant derivative}

Let $\tau$ be a $\Ck{2}$ map between open subsets of locally convex spaces. We write
$\dd\tau(y)(a)$ for its derivative and $\dd^{(2)}\tau(y; a, b)$ for the derivative of
$y \mapsto \dd\tau(y)(b)$ in the direction $a$; by the symmetry of second derivatives the latter is
symmetric in $(a, b)$. For two charts $i, j$ of the atlas of $M$ we write $\tau = j \circ i^{-1}$
for the transition map, defined on $i(U_i \cap U_j)$.

\begin{definition}[Connection]
  \label{def:connection}
  \lean{LCManifolds.Connection}
  \leanok
  \uses{def:mb-manifold, def:mb-groupoid}
  A \emph{connection} on $M$ assigns to each chart $i$ of the atlas a family of maps
  $\Gamma_i(x) \colon E \times E \to E$, $x \in i(U_i)$, such that
  \begin{enumerate}
    \item each $\Gamma_i(x)$ is bilinear, and $(x, a, b) \mapsto \Gamma_i(x)(a, b)$ is smooth on
    $i(U_i) \times E \times E$;
    \item (change of charts law) for every pair of charts $i, j$, every $y$ in the domain of
    $\tau = j \circ i^{-1}$ and all $a, b \in E$,
    \[
      \Gamma_j(\tau y)\bigl(\dd\tau(y)(a), \dd\tau(y)(b)\bigr)
        = \dd\tau(y)\bigl(\Gamma_i(y)(a, b)\bigr) + \dd^{(2)}\tau(y; a, b).
    \]
  \end{enumerate}
  The maps $\Gamma_i$ are the \emph{Christoffel symbols} of the connection in the chart $i$.
\end{definition}

The symbols are not tensorial: they transform affinely, the defect being the second derivative of
the transition map. Our convention is that the first argument of $\Gamma$ is the direction in
which the base point moves. Continuity of each bilinear map $\Gamma_i(x)$ is a consequence of the
joint smoothness, and no topology on the space of bilinear maps is needed.

\begin{definition}[Local covariant derivative]
  \label{def:conn-local-cov-deriv}
  \lean{LCManifolds.localCovDeriv}
  \leanok
  Let $\Gamma$ be a family of bilinear maps on an open set $U \subseteq E$ and let
  $X, Y \colon U \to E$. The \emph{local covariant derivative} is
  \[
    (\nabla_X Y)(y) = \dd Y(y)\bigl(X(y)\bigr) - \Gamma(y)\bigl(X(y), Y(y)\bigr).
  \]
\end{definition}

\begin{lemma}[Change of coordinates]
  \label{lem:conn-local-transition}
  \lean{LCManifolds.dMB_comp_transition, LCManifolds.localCovDeriv_transition}
  \leanok
  \uses{def:conn-local-cov-deriv}
  Let $\tau$ be $\Ck{2}$, and let $X_i, Y_i$ and $X_j, Y_j$ be related by $\tau$, that is
  $X_j \circ \tau = \dd\tau \cdot X_i$ and $Y_j \circ \tau = \dd\tau \cdot Y_i$, with $Y_i, Y_j$
  of class $\Ck{1}$. If $\Gamma_i, \Gamma_j$ obey the change of charts law of
  Definition~\ref{def:connection}, then
  \[
    (\nabla_{X_j} Y_j)(\tau y) = \dd\tau(y)\bigl((\nabla_{X_i} Y_i)(y)\bigr).
  \]
\end{lemma}

\begin{proof}
  \leanok
  Differentiating $Y_j \circ \tau = \dd\tau \cdot Y_i$ in the direction $w$ gives, by the
  second-order chain rule,
  $\dd Y_j(\tau y)(\dd\tau(y) w) = \dd^{(2)}\tau(y; w, Y_i y) + \dd\tau(y)(\dd Y_i(y) w)$. Take
  $w = X_i(y)$: the term $\dd^{(2)}\tau(y; X_i y, Y_i y)$ cancels with the one in the change of
  charts law of $\Gamma$, and the rest is linear in $\dd\tau(y)$.
\end{proof}

\begin{definition}[Covariant derivative]
  \label{def:conn-cov-deriv}
  \lean{LCManifolds.Connection.covDeriv, LCManifolds.Connection.isVectorField_covDeriv}
  \leanok
  \uses{def:connection, def:conn-local-cov-deriv, def:vector-field}
  Let $C$ be a connection on $M$ and $X, Y$ vector fields. The \emph{covariant derivative}
  $\nabla_X Y$ is the vector field whose value at $x$ is the local covariant derivative of the
  local representatives $\tilde X, \tilde Y$ of $X, Y$, computed in the preferred chart at $x$:
  $\dd\tilde Y \cdot \tilde X - \Gamma(\tilde X, \tilde Y)$.
\end{definition}

\begin{theorem}[Chart independence]
  \label{thm:conn-cov-deriv-chart}
  \lean{LCManifolds.Connection.localRep_covDeriv}
  \leanok
  \uses{def:conn-cov-deriv, lem:conn-local-transition, def:tangent-bundle}
  Let $Y$ be a $\Ck{1}$ vector field and $X$ any vector field. In every chart $j$ of the atlas,
  the local representative of $\nabla_X Y$ is
  $\dd\tilde Y_j \cdot \tilde X_j - \Gamma_j(\tilde X_j, \tilde Y_j)$.
\end{theorem}

\begin{proof}
  \leanok
  Apply Lemma~\ref{lem:conn-local-transition} to the transition map from the preferred chart to
  $j$, which is smooth; the local representatives of vector fields are related by its
  derivative.
\end{proof}

So the choice of the preferred chart in Definition~\ref{def:conn-cov-deriv} is only a matter of
convenience. The usual algebraic rules follow at once from the local formula.

\begin{proposition}[Rules of the covariant derivative]
  \label{prop:conn-rules}
  \lean{LCManifolds.Connection.covDeriv_smulField_left, LCManifolds.Connection.covDeriv_smulField_right, LCManifolds.Connection.contMDiffMB_covDeriv}
  \leanok
  \uses{thm:conn-cov-deriv-chart, def:contmdiff}
  Let $X, Y$ be vector fields with $Y$ of class $\Ck{1}$, and let $\varphi \colon M \to \R$.
  \begin{enumerate}
    \item $\nabla_{\varphi X} Y = \varphi\, \nabla_X Y$.
    \item If $\varphi$ is $\Ck{1}$, then $\nabla_X(\varphi Y) = X(\varphi)\, Y + \varphi\,\nabla_X Y$
    (Leibniz rule).
    \item If $X$ is $\Ck{k}$ and $Y$ is $\Ck{k+1}$, then $\nabla_X Y$ is $\Ck{k}$.
  \end{enumerate}
\end{proposition}

\begin{proof}
  \leanok
  The first two items are computed in the preferred chart: $\Gamma$ is linear in its first
  argument, and the derivative of $\varphi \tilde Y$ in the direction $\tilde X$ is
  $\dd\varphi(\tilde X)\,\tilde Y + \varphi\,\dd\tilde Y(\tilde X)$. For the third, in each chart
  $y \mapsto \dd\tilde Y(y)(\tilde X y)$ and $y \mapsto \Gamma(y)(\tilde X y, \tilde Y y)$ are
  compositions of $\Ck{k}$ maps.
\end{proof}

\begin{proposition}[Torsion]
  \label{prop:conn-torsion}
  \lean{LCManifolds.Connection.covDeriv_sub_covDeriv_sub_lieBracket, LCManifolds.Connection.IsSymmetric, LCManifolds.Connection.covDeriv_sub_covDeriv}
  \leanok
  \uses{def:conn-cov-deriv}
  For vector fields $X, Y$, the vector $\nabla_X Y - \nabla_Y X - [X, Y]$ at $x$ is, in the
  preferred chart, $\Gamma(\tilde Y, \tilde X) - \Gamma(\tilde X, \tilde Y)$; in particular it only
  depends on $X(x)$ and $Y(x)$. We say that the connection is \emph{symmetric} (torsion free) if
  its symbols are symmetric in every chart; in that case $\nabla_X Y - \nabla_Y X = [X, Y]$.
\end{proposition}

\begin{proof}
  \leanok
  The Lie bracket in a chart is $\dd\tilde Y \cdot \tilde X - \dd\tilde X \cdot \tilde Y$, and
  the derivative terms cancel.
\end{proof}

\begin{remark}
  Symmetry of a connection is a chart-free condition: the change of charts law shows that if
  $\Gamma_i$ is symmetric then so is $\Gamma_j$, because $\dd^{(2)}\tau$ is symmetric and
  $\dd\tau(y)$ is invertible on the overlap. In the Lean code the definition quantifies over all
  the charts of the atlas, which avoids this argument.
\end{remark}

\section{Curvature}

We use Lang's sign for the curvature.

\begin{definition}[Curvature]
  \label{def:conn-curvature}
  \lean{LCManifolds.dSymb, LCManifolds.localCurv, LCManifolds.Connection.curvature}
  \leanok
  \uses{def:conn-cov-deriv}
  For vector fields $X, Y, Z$, the \emph{curvature} at $x$ is
  \[
    R(X, Y) Z = \nabla_X \nabla_Y Z - \nabla_Y \nabla_X Z - \nabla_{[X, Y]} Z .
  \]
  Locally, for a family of symbols $\Gamma$ on an open set of $E$, we write
  $\dd\Gamma_x(v)(a, b)$ for the derivative of $x \mapsto \Gamma(x)(a, b)$ in the direction $v$,
  and
  \[
    R_x(u, v) w = \dd\Gamma_x(v)(u, w) - \dd\Gamma_x(u)(v, w)
      + \Gamma_x\bigl(u, \Gamma_x(v, w)\bigr) - \Gamma_x\bigl(v, \Gamma_x(u, w)\bigr).
  \]
\end{definition}

\begin{lemma}[Local curvature formula]
  \label{lem:conn-local-curv}
  \lean{LCManifolds.dMB_symb_comp, LCManifolds.localCovDeriv_curv}
  \leanok
  \uses{def:conn-curvature, def:conn-local-cov-deriv}
  Let $X, Y, Z$ be $\Ck{2}$ fields on an open set $U \subseteq E$ and $\Gamma$ a $\Ck{1}$ family of
  bilinear maps. Then, pointwise on $U$,
  $\nabla_X \nabla_Y Z - \nabla_Y \nabla_X Z - \nabla_{[X, Y]} Z = R(X, Y) Z$
  with $R$ the trilinear expression of Definition~\ref{def:conn-curvature}.
\end{lemma}

\begin{proof}
  \leanok
  Expand both covariant derivatives. The derivative of $y \mapsto \Gamma(y)(A y, B y)$ is
  $\dd\Gamma_y(h)(A y, B y) + \Gamma_y(\dd A(y) h, B y) + \Gamma_y(A y, \dd B(y) h)$. All the
  terms containing second derivatives of $Z$ cancel by the symmetry of second derivatives, and
  the terms containing first derivatives of $X$, $Y$, $Z$ cancel with $\nabla_{[X,Y]} Z$ and
  among themselves. This is Lang \cite{Lan99}, Prop.~IX.1.2.
\end{proof}

\begin{theorem}[The curvature is tensorial]
  \label{thm:conn-curvature-tensorial}
  \lean{LCManifolds.Connection.curvature_eq}
  \leanok
  \uses{lem:conn-local-curv, thm:conn-cov-deriv-chart}
  Let $X, Y, Z$ be $\Ck{2}$ vector fields. Then $R(X, Y) Z$ at $x$ is the trilinear expression
  $R_{\tilde x}(\tilde X, \tilde Y)\tilde Z$ in the preferred chart, evaluated at the values of
  $X, Y, Z$ at $x$. In particular, it only depends on $X(x)$, $Y(x)$ and $Z(x)$.
\end{theorem}

\begin{proof}
  \leanok
  Near $x$ the local representatives of $\nabla_Y Z$, $\nabla_X Z$ and $[X, Y]$ in the preferred
  chart at $x$ are given by the local formulas (Theorem~\ref{thm:conn-cov-deriv-chart}), and the
  local covariant derivative at a point only depends on the germ of the second field. Then apply
  Lemma~\ref{lem:conn-local-curv}. Compare \cite{Lan99}, Prop.~IX.1.1--1.2.
\end{proof}

\begin{proposition}[Symmetries of the curvature]
  \label{prop:conn-curvature-symmetries}
  \lean{LCManifolds.localCurv_swap, LCManifolds.localCurv_bianchi, LCManifolds.Connection.curvature_swap, LCManifolds.Connection.curvature_bianchi}
  \leanok
  \uses{thm:conn-curvature-tensorial, prop:conn-torsion}
  For $\Ck{2}$ vector fields, $R(Y, X) Z = -R(X, Y) Z$. If the connection is symmetric, the first
  Bianchi identity holds: $R(X, Y) Z + R(Y, Z) X + R(Z, X) Y = 0$.
\end{proposition}

\begin{proof}
  \leanok
  Both identities are checked on the trilinear expression. Skew-symmetry is immediate. For
  Bianchi, if $\Gamma$ is symmetric then so is $\dd\Gamma_x(v)$, and the twelve terms of the
  cyclic sum cancel in pairs.
\end{proof}

\section{Sprays}

A spray, or second order vector field, is a vector field $\xi$ on $TM$ such that
$T\pi \circ \xi = \id$. In a chart it has the form $(x, v) \mapsto (x, v; v, f(x, v))$, and it is
determined by its \emph{local part} $f$. As we did for connections, we encode sprays through their
local parts.

\begin{definition}[Spray]
  \label{def:spray}
  \lean{LCManifolds.Spray, LCManifolds.Spray.IsQuadratic}
  \leanok
  \uses{def:mb-manifold, def:tangent-bundle}
  A \emph{spray} on $M$ assigns to each chart $i$ of the atlas a smooth map
  $f_i \colon i(U_i) \times E \to E$ such that, for every pair of charts $i, j$, every $y$ in the
  domain of $\tau = j \circ i^{-1}$ and every $v \in E$,
  \[
    f_j\bigl(\tau y, \dd\tau(y)(v)\bigr) = \dd\tau(y)\bigl(f_i(y, v)\bigr) + \dd^{(2)}\tau(y; v, v).
  \]
  The spray is \emph{quadratic} if $f_i(x, s v) = s^2 f_i(x, v)$ for all $i$, $x$, $v$ and
  $s \in \R$.
\end{definition}

The change of charts law is the change of charts law of the double tangent bundle $T(TM)$,
restricted to the points of the form $(x, v; v, w)$. In \cite{Lan99}, \S IV.3, a spray is by
definition quadratic; we keep the two notions apart, since the quadratic condition is what is
needed to pass to connections.

\begin{proposition}[The spray of a connection]
  \label{prop:conn-to-spray}
  \lean{LCManifolds.Connection.toSpray, LCManifolds.Connection.toSpray_isQuadratic}
  \leanok
  \uses{def:connection, def:spray}
  Every connection defines a quadratic spray, with local parts $f_i(x, v) = \Gamma_i(x)(v, v)$.
\end{proposition}

\begin{proof}
  \leanok
  Restrict the change of charts law of $\Gamma$ to the diagonal $a = b = v$. Homogeneity of degree
  two follows from bilinearity.
\end{proof}

The converse requires recovering a bilinear map from its quadratic form, and checking that the
polarization obeys the change of charts law of connections. For the first point, the following
identity holds without any regularity assumption.

\begin{lemma}[Homogeneous maps of degree two]
  \label{lem:conn-homogeneous}
  \lean{LCManifolds.hasDirDerivAt_of_homogeneous, LCManifolds.dMB_dMB_zero_of_homogeneous, LCManifolds.polarization}
  \leanok
  Let $g \colon E \to F$ satisfy $g(s v) = s^2 g(v)$ for all $s \in \R$, $v \in E$. Then
  $\dd^{(2)} g(0; v, v) = 2\, g(v)$. Moreover, if $B$ is symmetric and bilinear, then
  $B(a + b, a + b) - B(a - b, a - b) = 4 B(a, b)$.
\end{lemma}

\begin{proof}
  \leanok
  By homogeneity, $g(s v + t v) - g(s v) = \bigl((s + t)^2 - s^2\bigr) g(v)$, so the difference
  quotient can be computed exactly and $\dd g(s v)(v) = 2 s\, g(v)$ for every $s$.
  Differentiating once more along the line $s \mapsto s v$ at $s = 0$ gives $2 g(v)$. The
  polarization identity is immediate.
\end{proof}

\begin{theorem}[The connection of a quadratic spray]
  \label{thm:conn-spray-to-connection}
  \lean{LCManifolds.Spray.polar, LCManifolds.Spray.polar_eq, LCManifolds.Spray.polar_transition, LCManifolds.Spray.toConnection, LCManifolds.Spray.toConnection_isSymmetric, LCManifolds.Spray.toConnection_toSpray}
  \leanok
  \uses{def:spray, def:connection, lem:conn-homogeneous, prop:conn-to-spray}
  Let $f$ be a quadratic spray. In each chart $i$ put
  \[
    \Gamma_i(x)(a, b) = \tfrac14\bigl(f_i(x, a + b) - f_i(x, a - b)\bigr)
      = \tfrac12\, \dd^{(2)}\bigl(f_i(x, \cdot)\bigr)(0; a, b).
  \]
  Then the $\Gamma_i$ are the symbols of a symmetric connection, whose spray is $f$.
\end{theorem}

\begin{proof}
  \leanok
  By Lemma~\ref{lem:conn-homogeneous} applied to $g = f_i(x, \cdot)$ and to the symmetric
  bilinear map $\tfrac12 \dd^{(2)} g(0; \cdot, \cdot)$, the two expressions agree; the second one
  is bilinear and symmetric, and $\Gamma_i(x)(v, v) = f_i(x, v)$. Smoothness of the symbols
  follows from the first expression. The change of charts law for $\Gamma$ is obtained by
  polarizing the change of charts law of the spray, using that $\dd\tau(y)$ is linear and
  $\dd^{(2)}\tau(y; \cdot, \cdot)$ is symmetric and bilinear. This is \cite{Lan99}, Prop.~IV.3.2
  and \S VIII.2.
\end{proof}

Thus quadratic sprays and symmetric connections are the same object, through
$f(x, v) = \Gamma_x(v, v)$; in the rest of the chapter we work with connections.

\section{Covariant derivative along curves and geodesics}

We use the $\Ck{k}$ curves of Definition~\ref{def:ck-curve} and the velocity of
Definition~\ref{def:velocity}. Let $\alpha \colon \R \to M$ be a curve. A \emph{lift} of $\alpha$
is a map $\beta \colon \R \to TM$ with $\beta(t) \in T_{\alpha(t)} M$. In a chart $i$ we write
$\tilde\alpha = i \circ \alpha$ and $\mu_i(t)$ for the fiber coordinate of $\beta(t)$.

\begin{definition}[Lifts of curves]
  \label{def:conn-lift}
  \lean{LCManifolds.liftCoord, LCManifolds.IsCkLift}
  \leanok
  \uses{def:ck-curve, def:tangent-bundle}
  A lift $\beta$ of $\alpha$ is $\Ck{k}$ if, for every chart $i$, its fiber coordinate $\mu_i$ is
  $\Ck{k}$ on $\alpha^{-1}(U_i)$.
\end{definition}

\begin{lemma}[Change of charts along a curve]
  \label{lem:conn-curve-transition}
  \lean{LCManifolds.tangentChart_snd_transition, LCManifolds.dMB_dMB_curve}
  \leanok
  \uses{def:conn-lift, def:velocity}
  For two charts $i, j$ and $\tau = j \circ i^{-1}$, we have
  $\mu_j(t) = \dd\tau(\tilde\alpha_i t)(\mu_i t)$ and
  $\dot{\tilde\alpha}_j = \dd\tau(\tilde\alpha_i)(\dot{\tilde\alpha}_i)$, and
  \[
    \dot\mu_j = \dd^{(2)}\tau\bigl(\tilde\alpha_i; \dot{\tilde\alpha}_i, \mu_i\bigr)
      + \dd\tau(\tilde\alpha_i)(\dot\mu_i).
  \]
  The fiber coordinate of the velocity in any chart $j$ is $\dot{\tilde\alpha}_j$, and the
  velocity of a $\Ck{k+1}$ curve is a $\Ck{k}$ lift.
\end{lemma}

\begin{proof}
  \leanok
  The first identity is the change of charts of the tangent bundle, and the others follow from
  the chain rule and the product rule applied to $t \mapsto \dd\tau(\tilde\alpha_i t)(\mu_i t)$.
\end{proof}

\begin{definition}[Covariant derivative along a curve]
  \label{def:cov-deriv-along}
  \lean{LCManifolds.Connection.covDerivAlong}
  \leanok
  \uses{def:connection, def:conn-lift}
  Let $C$ be a connection and $\beta$ a lift of $\alpha$. The \emph{covariant derivative of
  $\beta$ along $\alpha$} at $t$ is, in the preferred chart at $\alpha(t)$,
  \[
    D_t \beta = \dot\mu - \Gamma(\tilde\alpha)\bigl(\dot{\tilde\alpha}, \mu\bigr).
  \]
\end{definition}

\begin{theorem}[Chart independence along curves]
  \label{thm:conn-cov-deriv-along-chart}
  \lean{LCManifolds.Connection.covDerivAlong_chart}
  \leanok
  \uses{def:cov-deriv-along, lem:conn-curve-transition}
  Let $\alpha$ be a $\Ck{1}$ curve and $\beta$ a $\Ck{1}$ lift. Read in any chart $j$ around
  $\alpha(t)$, the covariant derivative $D_t\beta$ is
  $\dot\mu_j - \Gamma_j(\tilde\alpha_j)(\dot{\tilde\alpha}_j, \mu_j)$.
\end{theorem}

\begin{proof}
  \leanok
  By Lemma~\ref{lem:conn-curve-transition}, $\dot\mu_j$ contains the term
  $\dd^{(2)}\tau(\tilde\alpha_i; \dot{\tilde\alpha}_i, \mu_i)$, which cancels with the one in the
  change of charts law of $\Gamma$ applied to $(\dot{\tilde\alpha}_i, \mu_i)$. The remaining
  terms are $\dd\tau(\tilde\alpha_i)$ applied to the expression in the chart $i$. Compare
  \cite{Lan99}, \S VIII.3.
\end{proof}

\begin{definition}[Geodesic]
  \label{def:geodesic}
  \lean{LCManifolds.Connection.IsGeodesic}
  \leanok
  \uses{def:cov-deriv-along, def:velocity}
  A curve $\alpha$ is a \emph{geodesic} of the connection $C$ if its velocity is parallel:
  $D_t \alpha' = 0$ for every $t$.
\end{definition}

\begin{theorem}[The geodesic equation]
  \label{thm:conn-geodesic-equation}
  \lean{LCManifolds.Connection.isGeodesic_iff}
  \leanok
  \uses{def:geodesic, thm:conn-cov-deriv-along-chart}
  A $\Ck{2}$ curve $\alpha$ is a geodesic if and only if, in every chart $j$,
  \[
    \ddot{\tilde\alpha}_j = \Gamma_j(\tilde\alpha_j)\bigl(\dot{\tilde\alpha}_j, \dot{\tilde\alpha}_j\bigr).
  \]
\end{theorem}

\begin{proof}
  \leanok
  By Theorem~\ref{thm:conn-cov-deriv-along-chart} and Lemma~\ref{lem:conn-curve-transition}, in
  the chart $j$ the covariant derivative of the velocity is
  $\ddot{\tilde\alpha}_j - \Gamma_j(\dot{\tilde\alpha}_j, \dot{\tilde\alpha}_j)$, transported by
  the invertible map $\dd\tau$. For the converse it suffices to use the preferred chart.
\end{proof}

In terms of the spray of the connection (Proposition~\ref{prop:conn-to-spray}), the geodesic
equation reads $\ddot{\tilde\alpha} = f(\tilde\alpha, \dot{\tilde\alpha})$ in every chart, which is
the classical definition of the geodesics of a spray.

\begin{remark}[Existence of geodesics]
  \label{rem:conn-no-existence}
  We do not claim that geodesics exist, nor that they are unique, for given initial position and
  velocity. In a Banach space the geodesic equation is an ordinary differential equation with a
  smooth right hand side, and the Picard--Lindel\"of theorem gives local existence and
  uniqueness; this is how the exponential map is constructed in \cite{Lan99} and \cite{Lar22}.
  In a general locally convex space, and even in Fr\'echet spaces, ordinary differential equations
  with smooth right hand side may have no solutions, or many \cite[Part~I, 5.6.1]{Ham82}. Therefore existence of geodesics,
  the exponential map and parallel transport along curves are not part of the theory developed
  here; they would have to be proved case by case, or under additional hypotheses on $E$.
  What we can do in general is to recognize geodesics, as in Theorem~\ref{thm:projector-geodesics}
  below. When the model space is a Banach space, existence and uniqueness are proved in the next
  section.
\end{remark}

\section{Geodesics on Banach manifolds}

In this section $E$ is a Banach space. Then Michal--Bastiani and Fr\'echet derivatives agree, and
in a chart $i$ the geodesic equation is the second order equation
\[
  c'' = \Gamma_i(c)(c', c')
\]
on the open set $\varphi_i(U_i) \subseteq E$, with $\Gamma_i$ smooth and bilinear in the last two
variables. We say that $\alpha$ is a \emph{geodesic on an open set} $I \subseteq \mathbb{R}$ if it is
continuous on $I$ and satisfies this equation at every $t \in I$ in the preferred chart at
$\alpha(t)$.

\begin{theorem}[The geodesic equation does not depend on the chart]
  \label{thm:geodesic-eq-chart}
  \lean{LCManifolds.geodesicEq_comp, LCManifolds.Connection.geodesicEqAt_of_chart}
  \leanok
  If $\alpha$ is continuous at $t$ and satisfies the geodesic equation at $t$ in a chart around
  $\alpha(t)$, it satisfies it in every chart around $\alpha(t)$.
\end{theorem}

\begin{proof}
  \leanok
  With $\tau = \varphi_j \circ \varphi_i^{-1}$ and $c = \varphi_i \circ \alpha$,
  $(\tau \circ c)'' = D^2\tau(c', c') + D\tau(c'')$, and the change of charts law of the symbols
  reads $\Gamma_j(\tau y)(D\tau\, a, D\tau\, b) = D\tau\, \Gamma_i(y)(a, b) + D^2\tau(a, b)$.
\end{proof}

\begin{theorem}[Local existence of geodesics, with uniform time]
  \label{thm:geodesic-existence}
  \lean{LCManifolds.Connection.exists_isGeodesicOn_uniform, LCManifolds.Connection.exists_isGeodesicOn}
  \leanok
  \uses{thm:geodesic-eq-chart}
  Let $y_0$ be a point of the chart $i$. There is $r > 0$ such that for every $K \ge 0$ there is
  $\varepsilon > 0$ with the following property: for every $y$ with $\|y - y_0\| \le r$ and every
  $w$ with $\|w\| \le K$ there is a geodesic $\alpha$ on $(-\varepsilon, \varepsilon)$, staying in
  the domain of $i$, with $\varphi_i(\alpha(0)) = y$ and $(\varphi_i \circ \alpha)'(0) = w$.
\end{theorem}

\begin{proof}
  \leanok
  Picard--Lindel\"of for the first order system $(c, v)' = (v, \Gamma_i(c)(v, v))$, whose right
  hand side is $\Ck{1}$, hence locally Lipschitz; the time is uniform for small velocities, and the
  homogeneity of the equation ($t \mapsto c(st)$ is a solution with initial velocity $s\, c'(0)$)
  makes it uniform for velocities bounded by $K$.
\end{proof}

\begin{theorem}[Uniqueness of geodesics]
  \label{thm:geodesic-uniqueness}
  \lean{LCManifolds.Connection.IsGeodesicOn.eqOn}
  \leanok
  \uses{thm:geodesic-eq-chart}
  If $M$ is Hausdorff, two geodesics on $(a, b)$ with the same position and velocity at some
  $t_0 \in (a, b)$ agree on $(a, b)$.
\end{theorem}

\begin{proof}
  \leanok
  The set of $t$ near which the two geodesics agree is open. It is also closed in $(a, b)$: at a
  limit point $t$ the two curves meet, because $M$ is Hausdorff; in a chart around this common
  point both solve the same equation, which has unique solutions since its right hand side is
  locally Lipschitz. By connectedness it is all of $(a, b)$.
\end{proof}

\begin{remark}
  The Hausdorff hypothesis cannot be dropped: on the line with two origins, the two copies of a
  line through the origins are geodesics of the flat connection with the same data at $t = 1$, and
  they differ at $t = 0$.
\end{remark}

\begin{theorem}[Extension of geodesics]
  \label{thm:geodesic-extension}
  \lean{LCManifolds.Connection.IsGeodesicOn.exists_extend}
  \leanok
  \uses{thm:geodesic-existence, thm:geodesic-uniqueness}
  Let $M$ be Hausdorff, and let $\alpha$ be a geodesic on $(a, b)$ such that $\alpha(t) \to q$ as
  $t \to b^-$ and the velocity of $\alpha$, read in a chart around $q$, stays bounded as
  $t \to b^-$. Then $\alpha$ extends to a geodesic on $(a, b')$ for some $b' > b$.
\end{theorem}

\begin{proof}
  \leanok
  Let $\varepsilon$ be the uniform time of Theorem~\ref{thm:geodesic-existence} around
  $\varphi_i(q)$. Take $t_1 > b - \varepsilon/2$ with $\alpha(t_1)$ close to $q$, and the geodesic
  $\mu$ with the position and velocity of $\alpha$ at $t_1$. By uniqueness
  $\mu(t - t_1) = \alpha(t)$ where both are defined, so the two curves glue to a geodesic on
  $(a, t_1 + \varepsilon)$, and $t_1 + \varepsilon > b$.
\end{proof}

\begin{theorem}[Complete Finsler manifolds: geodesics extend {\cite[Prop.~9.2.1]{Lar22}}]
  \label{thm:finsler-geodesic-extension}
  \lean{LCManifolds.Connection.IsGeodesicOn.exists_extend_of_complete}
  \leanok
  \uses{thm:geodesic-extension}
  Let $M$ be Hausdorff, with a Finsler length structure $\mu$ bounded below, such that $M$ is
  complete for the symmetrized distance $\max(d_b(x, y), d_b(y, x))$. If $\alpha$ is a geodesic on
  $(a, b)$ with $\mu(\dot\alpha)$ and $\mu(-\dot\alpha)$ bounded, then $\alpha$ extends to a
  geodesic on an interval $(a', b') \supsetneq (a, b)$ at both ends.
\end{theorem}

\begin{proof}
  \leanok
  Reparametrizing arcs of $\alpha$ by a smooth ramp gives $\Ck{1}$ competitors for $d_b$, so
  $d_b(\alpha(s), \alpha(t)) \le L(t - s)$ and the same backwards; thus $\alpha(t)$ is Cauchy as
  $t \to b^-$. Its limit in $(M, d_b)$ is a limit in $M$ because $\mu$ is bounded below, and the
  lower bound also bounds the velocity in a chart around the limit point.
  Theorem~\ref{thm:geodesic-extension} applies; the end $a$ is treated by reversing time.
\end{proof}

\begin{remark}
  In \cite{Lar22} the metric is invariant under the parallel transport of the spray, so geodesics
  have constant speed; we only use bounded speed. Since $\mu$ need not be reversible, both
  $\mu(\dot\alpha)$ and $\mu(-\dot\alpha)$ have to be bounded. Maximal geodesics are not
  constructed, so the statement ``every geodesic is defined on $\mathbb{R}$'' appears as the
  extension property above.
\end{remark}

\section{Fields of projectors: derive and project}
\label{sec:conn-projectors}

Let $F$ be a locally convex space and $X \subseteq F$ an embedded submanifold modeled on a closed
subspace $S \subseteq F$ (Definition~\ref{def:embedded-submanifold}). Recall that in a slice chart
$\varphi$ the tangent space at $p$ is $T_p X = \{w \in F : \dd\varphi(p)(w) \in S\}$, and that it
does not depend on the slice chart (Definition~\ref{def:tangent-space-submanifold}). We write
$\psi = \varphi^{-1}|_S$ for the local parametrization of $X$ given by $\varphi$; then
$\dd\psi(y)$ maps $S$ injectively onto $T_{\psi(y)} X$.

\begin{definition}[Field of projectors]
  \label{def:projector-field}
  \lean{LCManifolds.ProjectorField, LCManifolds.ProjectorField.param, LCManifolds.ProjectorField.localE}
  \leanok
  \uses{def:embedded-submanifold, def:tangent-space-submanifold}
  A \emph{field of projectors} onto the tangent spaces of $X$ is a map
  $E \colon X \times F \to F$, $(p, w) \mapsto E_p w$, such that
  \begin{enumerate}
    \item each $E_p$ is linear, $E_p^2 = E_p$, and the range of $E_p$ is $T_p X$;
    \item for every slice chart $\varphi$, the map $(y, w) \mapsto E_{\psi(y)} w$ is smooth on
    $\varphi(U) \cap S \times F$.
  \end{enumerate}
\end{definition}

\begin{remark}
  In \cite{Lar22} the field of projectors is a smooth map $X \to \mathcal B(F)$ into the bounded
  operators of a Banach space. For a locally convex space $F$ there is no good topology on the
  space of continuous operators, so we replace it by a map of two variables, linear in the second
  one, with joint smoothness. Note also that no complement of $S$ in $F$ is assumed: the kernels
  $\ker E_p$ are complements of the tangent spaces $T_p X$, but they are not required to come from
  a global splitting of $F$, nor to be related to each other.
\end{remark}

\begin{lemma}[Tangent vectors]
  \label{lem:projector-tangent}
  \lean{LCManifolds.ProjectorField.mem_tangentSpace, LCManifolds.ProjectorField.E_eq_self, LCManifolds.IsEmbeddedSubmanifold.dMB_symm_mem_tangentSpace, LCManifolds.ProjectorField.E_dMB_param, LCManifolds.ProjectorField.dMB_param_injective}
  \leanok
  \uses{def:projector-field}
  $E_p w \in T_p X$ for every $w$, and $E_p w = w$ if $w \in T_p X$. In a slice chart,
  $\dd\psi(y)(s) \in T_{\psi(y)} X$ for $s \in S$, hence $E_{\psi(y)}\dd\psi(y)(s) = \dd\psi(y)(s)$,
  and $\dd\psi(y)$ is injective on $S$.
\end{lemma}

\begin{proof}
  \leanok
  The first two claims follow from $E_p^2 = E_p$ and the description of the range. Since
  $\dd\varphi$ is a left and right inverse of $\dd\varphi^{-1}$, we have
  $\dd\varphi(\dd\psi(y) s) = s \in S$, which gives the rest.
\end{proof}

\begin{definition}[Symbols of a field of projectors]
  \label{def:projector-symbols}
  \lean{LCManifolds.ProjectorField.symb, LCManifolds.ProjectorField.symb_mem, LCManifolds.ProjectorField.coe_symb}
  \leanok
  \uses{def:projector-field}
  In a slice chart $\varphi$, with $\psi = \varphi^{-1}|_S$ and $y \in \varphi(U) \cap S$, put
  \[
    \Gamma_y(a, b) = -\,\dd\varphi(\psi y)\Bigl(E_{\psi(y)}\bigl(\dd^{(2)}\psi(y; a, b)\bigr)\Bigr),
    \qquad a, b \in S .
  \]
  This vector lies in $S$, because $E_{\psi(y)}(\cdot)$ is tangent.
\end{definition}

Equivalently, $\dd\psi(y)\bigl(\Gamma_y(a, b)\bigr) = -E\bigl(\dd^{(2)}\psi(y; a, b)\bigr)$: the
symbols are the tangent part of the second derivative of the parametrization, read in the chart.

\begin{lemma}[Properties of the symbols]
  \label{lem:projector-symbols}
  \lean{LCManifolds.ProjectorField.symb_symm, LCManifolds.ProjectorField.smoothMBOn_symb, LCManifolds.ProjectorField.dMB_param_symb, LCManifolds.ProjectorField.symb_transition}
  \leanok
  \uses{def:projector-symbols, lem:projector-tangent, def:connection}
  The symbols $\Gamma_y$ are bilinear and symmetric, they are smooth on the chart domain, and for
  two slice charts they obey the change of charts law of Definition~\ref{def:connection}.
\end{lemma}

\begin{proof}
  \leanok
  Symmetry and bilinearity come from those of $\dd^{(2)}\psi(y; \cdot, \cdot)$ (Schwarz), and
  smoothness from the smoothness of $E$ in slice charts and of $\psi$. For the change of charts
  law let $\tau$ be the transition map between two slice charts, so that $\psi' \circ \tau = \psi$
  near $y$. The second-order chain rule gives
  $\dd^{(2)}\psi = \dd^{(2)}\psi'(\dd\tau\, a, \dd\tau\, b) + \dd\psi'(\dd^{(2)}\tau(a, b))$. The
  last term is tangent, so $E$ fixes it; applying $-\dd\varphi' \circ E$ yields the law.
\end{proof}

\begin{theorem}[The connection of a field of projectors]
  \label{thm:projector-connection}
  \lean{LCManifolds.ProjectorField.connection, LCManifolds.ProjectorField.connection_isSymmetric}
  \leanok
  \uses{lem:projector-symbols, def:connection, prop:conn-torsion}
  A field of projectors $E$ onto the tangent spaces of $X$ defines a symmetric connection on $X$,
  whose symbols in the chart induced by a slice chart $\varphi$ are
  $\Gamma_y(a, b) = -\dd\varphi\bigl(E(\dd^{(2)}\psi(y; a, b))\bigr)$.
\end{theorem}

\begin{proof}
  \leanok
  The charts of the atlas of $X$ are the restrictions of slice charts, and
  Lemma~\ref{lem:projector-symbols} provides all the axioms. This is the first part of
  \cite[Teorema~6.6.1]{Lar22}.
\end{proof}

The next result explains the title of the section: the covariant derivative is obtained by
differentiating in $F$ and projecting onto the tangent space.

\begin{theorem}[Derive and project]
  \label{thm:projector-derive-project}
  \lean{LCManifolds.ProjectorField.dMB_param_localCovDeriv}
  \leanok
  \uses{thm:projector-connection, def:conn-local-cov-deriv}
  Let $A, B$ be fields on an open subset of the chart domain in $S$, with $B$ of class $\Ck{1}$,
  and let $\hat B(y) = \dd\psi(y)(B y)$ be the corresponding tangent field read in $F$. Then
  \[
    \dd\psi(y)\bigl((\nabla_A B)(y)\bigr) = E_{\psi(y)}\bigl(\dd\hat B(y)(A y)\bigr).
  \]
  In other words, $\nabla_X Y = E(\dd Y \cdot X)$ when the fields are read in the ambient
  space $F$.
\end{theorem}

\begin{proof}
  \leanok
  By the second-order chain rule,
  $\dd\hat B(y)(A y) = \dd^{(2)}\psi(y; A y, B y) + \dd\psi(y)(\dd B(y)(A y))$. Apply $E$: the
  second term is tangent and is fixed, and $E(\dd^{(2)}\psi(y; A y, B y))
  = -\dd\psi(y)(\Gamma_y(A y, B y))$ by Definition~\ref{def:projector-symbols}. Compare
  \cite[\S6.6]{Lar22}.
\end{proof}

\begin{lemma}[Acceleration in a slice chart]
  \label{lem:projector-accel}
  \lean{LCManifolds.ProjectorField.E_accel_eq, LCManifolds.IsEmbeddedSubmanifold.isCkCurve_of_smooth}
  \leanok
  \uses{def:projector-symbols, def:ck-curve}
  Let $c$ be a $\Ck{2}$ curve in the chart domain in $S$ and $\alpha = \psi \circ c$. Then
  \[
    E_{\alpha(t)}\bigl(\ddot\alpha(t)\bigr)
      = \dd\psi(c t)\bigl(\ddot c(t) - \Gamma_{c(t)}(\dot c(t), \dot c(t))\bigr),
  \]
  where $\ddot\alpha$ is the acceleration of $\alpha$ in $F$. Moreover, a curve in $X$ which is
  smooth as a curve in $F$ is a $\Ck{k}$ curve of $X$ for every $k$.
\end{lemma}

\begin{proof}
  \leanok
  We have $\dot\alpha = \dd\psi(c)(\dot c)$ and, by the chain rule and the product rule,
  $\ddot\alpha = \dd^{(2)}\psi(c; \dot c, \dot c) + \dd\psi(c)(\ddot c)$; apply $E$ as in
  the proof of Theorem~\ref{thm:projector-derive-project}. For the last claim, the local
  representative in a slice chart is the corestriction to $S$ of $\varphi \circ \alpha$.
\end{proof}

\begin{theorem}[Derive and project: geodesics]
  \label{thm:projector-geodesics}
  \lean{LCManifolds.ProjectorField.connection_isGeodesic_iff}
  \leanok
  \uses{thm:projector-connection, thm:conn-geodesic-equation, lem:projector-accel, def:geodesic}
  Let $\alpha$ be a $\Ck{2}$ curve in $X$. Then $\alpha$ is a geodesic of the connection defined
  by the field of projectors $E$ if and only if its acceleration in $F$ is normal:
  \[
    E_{\alpha(t)}\bigl(\ddot\alpha(t)\bigr) = 0 \qquad \text{for every } t.
  \]
\end{theorem}

\begin{proof}
  \leanok
  By Theorem~\ref{thm:conn-geodesic-equation}, $\alpha$ is a geodesic if and only if
  $\ddot c = \Gamma_c(\dot c, \dot c)$ in every chart, where $c$ is the local representative. By
  Lemma~\ref{lem:projector-accel} this is equivalent to $E(\ddot\alpha) = 0$, since $\dd\psi(c)$ is
  injective on $S$ (Lemma~\ref{lem:projector-tangent}). This is part of \cite[Teorema~6.6.1]{Lar22}.
\end{proof}

For instance, for the unit sphere of a (complete) Hilbert space with the orthogonal projectors
onto $p^\perp$, the geodesics are the curves in the sphere with $\ddot\alpha(t)$ proportional to
$\alpha(t)$. This example is not formalized.

Finally, we compare with the formulation of the book, where the connection is written in terms of
the derivative of the field of projectors. We denote with $E_{*p, v} w$ the derivative of
$q \mapsto E_q w$ at $p$ in the direction $v$, computed in the slice chart.

\begin{proposition}[The identity of the book]
  \label{prop:projector-book-identity}
  \lean{LCManifolds.ProjectorField.dMB_localE_param, LCManifolds.ProjectorField.dMB_localE_param_symm}
  \leanok
  \uses{def:projector-field, lem:projector-tangent}
  Let $v = \dd\psi(y)(a)$ and $w = \dd\psi(y)(b)$ be tangent vectors at $p = \psi(y)$, with
  $a, b \in S$. Then
  \[
    E_{*p, v}\, w = (1 - E_p)\, \dd^{(2)}\psi(y; a, b).
  \]
  In particular $E_{*p, v}\, w = E_{*p, w}\, v$.
\end{proposition}

\begin{proof}
  \leanok
  Since $\dd\psi(y')(b)$ is tangent, $E_{\psi(y')}\dd\psi(y')(b) = \dd\psi(y')(b)$ for $y'$ near
  $y$. Differentiate this identity in the direction $a$: by linearity of $E$ in the second
  variable, the left side gives $E_{*p, v} w + E_p\, \dd^{(2)}\psi(y; a, b)$, and the right side
  gives $\dd^{(2)}\psi(y; a, b)$. The symmetry follows from the symmetry of $\dd^{(2)}\psi$.
\end{proof}

\begin{remark}
  In \cite{Lar22} the connection is presented through $\Gamma_p(v, w) = E_{*p, v} w$, which by
  Proposition~\ref{prop:projector-book-identity} is the normal part $(1 - E)\dd^{(2)}\psi(a, b)$
  of the second derivative, read in $F$. The symbols of
  Theorem~\ref{thm:projector-connection} are instead the tangent part, pulled back to $S$ by the
  chart; the two descriptions are related through the decomposition
  $\dd^{(2)}\psi = E\,\dd^{(2)}\psi + (1 - E)\dd^{(2)}\psi$. In both cases the symmetry of the
  connection is a consequence of the symmetry of second derivatives of $\psi$. In the Lean code
  the connection is defined by the tangent part, and the book's formula enters only through
  Proposition~\ref{prop:projector-book-identity}.
\end{remark}
